\section{Complexity and practical status}
\label{sec:automated-reasoning-07}

A \emph{decision problem} is, under the fixed coding, a set of natural
numbers.  It is \emph{polynomial-time decidable} iff some deterministic
Turing machine decides it within \(p(n)\) steps on every input of length
\(n\), for some polynomial \(p\).  The class \(\NP\) consists of the
decision problems \(A\) for which there are a polynomial-time decidable
binary relation \(R\) and a polynomial \(q\) such that \(x\in A\) iff
some witness \(w\) whose length is at most \(q(n)\), where \(n\) is
the length of \(x\), satisfies \(R(x,w)\).  A problem is \(\NP\)-\emph{complete} iff
it belongs to \(\NP\) and every member of \(\NP\) reduces to it by a
polynomial-time computable many-one reduction.  Propositional
satisfiability is \(\NP\)-complete \cite{Cook1971}.  Whether
propositional satisfiability is decidable in worst-case polynomial time is
the open \(\mathsf P\) versus \(\NP\) question; no worst-case
polynomial-time decision procedure is known.  This complexity
result does not imply that every large instance is difficult; modern solvers
exploit structure.  It also does not justify replacing a proof with a solver
response whose trust path is unstated.

First-order validity is computably enumerable and undecidable.  A complete
first-order prover may therefore search forever on some nonvalid inputs.  A
timeout means only that the configured run did not finish.

\begin{method}[Automated-result audit]
Record:
\begin{enumerate}
  \item the exact input language and version;
  \item the intended logical semantics;
  \item every declaration and nonlogical axiom;
  \item transformations, including whether they preserve equivalence or only
        satisfiability;
  \item the calculus and search program;
  \item the result status and resource limit;
  \item the proof certificate or model, when available; and
  \item the independent checker and its trusted base.
\end{enumerate}
\end{method}
